% Источники: exam/materials/моделирование.pdf, раздел 3.
\section{Натурный, полунатурный, вычислительный эксперименты.}

\textbf{Натурный эксперимент} — эксперимент с реальными объектами и процессами.

\textbf{Пример:} исследование аэродинамических свойств автомобиля. В этом случае автомобиль помещается в аэродинамическую трубу, где измеряется сопротивление воздуха при различных скоростях.

\textbf{Достоинства:} высокая точность решения задачи.

\textbf{Недостатки:} такие эксперименты дорогостоящи.

\textbf{Полунатурный эксперимент} — эксперимент, где в качестве модели используются физическая модель объекта, доступные характеристики оцениваются экспериментально, другие — программно.

\textbf{Пример:} испытание систем управления самолета. Модель самолета (или его крыла) создается в аэродинамической трубе. Здесь измеряются силы сопротивления воздуха, подъемная сила и другие физические параметры. Однако поведение самолета в реальных условиях полета (например, при изменении высоты, температуры воздуха или ветра) моделируется на компьютере.

\textbf{Достоинства:} дешевле натурного.

\textbf{Недостатки:} снижается точность по сравнению с натурным. Требования к точности эксперимента определяются целью.

\textbf{Вычислительный эксперимент} — эксперимент, в основе которого лежит реализация математическом модели на ЦВМ.

\textbf{Пример:} моделирование распространения эпидемии.

\textbf{Достоинства:} позволяет изменить время протекания эксперимента, установить зависимости между наблюдаемыми показателями и исследуемыми.

\textbf{Недостатки:} точность результатов решения задачи ниже точности результатов натурного эксперимента.

\clearpage
